\section{Background}
\label{sec:background}

\subsection{State Estimation: Filtering and Smoothing}
\label{subsec:state-estimation}

State estimation infers the motion of an object from noisy observations.
For planar vehicle motion, a common state contains position and velocity,
$\mathbf{x}_k=(p_x,p_y,v_x,v_y)^{\mathsf T}$, at sampling instant $k$.
For Cartesian position and directional velocity measurements, a linear
state-space model describes its evolution and observation as
\begin{equation}
  \mathbf{x}_k=\mathbf{F}_k\mathbf{x}_{k-1}+\mathbf{w}_k,
  \qquad
  \mathbf{z}_k=\mathbf{H}_k\mathbf{x}_k+\mathbf{v}_k,
  \label{eq:state-space}
\end{equation}
where $\mathbf{F}_k$ models motion, $\mathbf{H}_k$ relates the state to the
measurements, and $\mathbf{w}_k$ and $\mathbf{v}_k$ represent process and
measurement noise, with covariances $\mathbf{Q}_k$ and $\mathbf{R}_k$.
In a constant-velocity model, position advances by velocity times the
sampling interval; process noise accounts for departures from that
prediction, including acceleration~\cite{barshalom2001estimation}.

The Kalman filter recursively propagates a state estimate and its error
covariance, then updates them using the next measurement and its
uncertainty~\cite{kalman1960filtering}. Filtering estimates the state at
$k$ using observations through $k$. Fixed-interval smoothing estimates
that same state using the entire observation sequence, including later
measurements. The Rauch--Tung--Striebel (RTS) smoother performs a backward
recursion over the filtered estimates~\cite{rauch1965smoothing}.
This distinction is relevant to retrospective analysis, where the complete
recording is available. Both estimates depend on the assumed motion model;
an abrupt velocity change at impact requires that assumption to be revisited.

\subsection{Multi-Object Tracking and Data Association}
\label{subsec:multi-object-tracking}

Multi-object tracking estimates several trajectories while maintaining
their identities over time. Sensor detections generally carry no reliable
object identity: a vehicle can produce several returns, disappear from
view, or pass close to another vehicle. Data association determines which
observations belong to which tracks, while track management handles their
initiation, continuation, and termination~\cite{barshalom2001estimation}.
For a predicted state $\hat{\mathbf{x}}$ with covariance $\mathbf{P}$,
the measurement innovation $\mathbf{r}=\mathbf{z}-\mathbf{H}\hat{\mathbf{x}}$
has covariance $\mathbf{S}=\mathbf{H}\mathbf{P}\mathbf{H}^{\mathsf T}+\mathbf{R}$.
Compatibility can be assessed using the squared Mahalanobis distance,
\begin{equation}
  d^2=\mathbf{r}^{\mathsf T}\mathbf{S}^{-1}\mathbf{r}.
  \label{eq:association-distance}
\end{equation}
A statistical gate rejects incompatible candidates. Passing the gate
establishes compatibility, but several tracks may still explain the same
observation.

Automotive radar also requires an extended-object view of a vehicle.
Multiple detections can arise from spatially separated reflection centres,
and their distribution depends on the visible parts of the
body~\cite{hammarstrand2012extended}. Doppler measures relative velocity along each
return's line of sight; a single return does not directly determine the
tangential component. Consequently, a tracked radar position need not
coincide with the vehicle centre, and different observers can track
different surfaces of the same car. Extended-object tracking explicitly
accounts for the relationship between object motion, physical extent,
and multiple detections~\cite{granstrom2017extended}.

With several observers, association also operates between tracks.
A local track label identifies an observation history within one source;
track-to-track association determines whether histories from different
sources describe the same physical object. Establishing that correspondence
and fusing the associated state estimates are distinct estimation
problems~\cite{chong2000association}.

\subsection{Traffic Conflict Measures}
\label{subsec:traffic-conflicts}

Traffic-conflict measures quantify proximity to a collision under specified
motion assumptions. Time to collision (TTC) is the time until the first
predicted contact if those assumptions continue to
hold~\cite{hayward1972near-miss}. For two vehicles following a straight path
at constant speeds, with follower speed $v_f>v_l$ and leader speed $v_l$,
TTC and the following time gap are
\begin{equation}
  \mathrm{TTC}=\frac{d}{v_f-v_l},
  \qquad
  \tau_{\mathrm{gap}}=\frac{d}{v_f},
  \label{eq:ttc-time-gap}
\end{equation}
where $d>0$ is the bumper-to-bumper clearance. The latter is a net time gap;
gross time headway instead uses front-to-front spacing. These measures
describe different properties: two vehicles travelling at the same speed
can have a short following gap while their constant-speed TTC is infinite.
Close following and an imminent collision are therefore distinct
conditions~\cite{vogel2003headway}.

For cut-ins and crossing traffic, the distance-to-closing-speed ratio alone
does not establish a collision course. Predictions must account for both
vehicles' paths and physical footprints. Measures based on braking further
consider the deceleration needed to avoid predicted contact after an
assumed response delay. Their interpretation depends on the motion model,
response assumptions, and available deceleration~\cite{westhofen2023criticality}.
They characterize aspects of conflict urgency without establishing its cause.

Response timing also depends on context. Naturalistic studies show that
braking onset varies with the evolving kinematic urgency of an encounter,
rather than following a universal delay after a single hazard
onset~\cite{markkula2016braking}. An interval from a conflict observable
in sensor data to brake application describes an observed braking delay;
it does not establish when the driver perceived the hazard or whether the
driver was distracted.
